\textbf{Phần 1. Trắc nghiệm lựa chọn}
\setcounterpageref{ex}{0}

\begin{ex} %[5D12V69]
Tính giá trị biểu thức bằng cách thuận tiện nhất: $4{,}85 \times 74 + 4{,}85 \times 25 + 4{,}85$. Kết quả là
\choice
{$48{,}5$}
{\True $485$}
{$4\,850$}
{$388$}
\loigiai{
Ta có:
$$4{,}85 \times 74 + 4{,}85 \times 25 + 4{,}85 \times 1 = 4{,}85 \times (74 + 25 + 1) = 4{,}85 \times 100 = 485.$$
\textbf{Đáp án của câu hỏi là B}
}
\end{ex}

\begin{ex} %[5D12V69]
Giá trị của biểu thức $\left(\dfrac{3}{4} + 1{,}25\right) \times \dfrac{2}{5} - 0{,}3$ là
\choice
{$0{,}4$}
{\True $0{,}5$}
{$0{,}6$}
{$0{,}8$}
\loigiai{
Đổi: $\dfrac{3}{4} = 0{,}75$.\\
Trong ngoặc: $0{,}75 + 1{,}25 = 2$.\\
Nhân với $\dfrac{2}{5} = 0{,}4$:
$$2 \times 0{,}4 = 0{,}8.$$
Trừ đi $0{,}3$:
$$0{,}8 - 0{,}3 = 0{,}5.$$
\textbf{Đáp án của câu hỏi là B}
}
\end{ex}

\begin{ex} %[5D12V69]
Tìm $x$ biết: $x \times 3{,}5 + x \times 6{,}5 = 42$.
\choice
{$4$}
{\True $4{,}2$}
{$4{,}5$}
{$5{,}2$}
\loigiai{
Ta có:
$$x \times (3{,}5 + 6{,}5) = 42$$
$$x \times 10 = 42 \Rightarrow x = 42 : 10 = 4{,}2.$$
\textbf{Đáp án của câu hỏi là B}
}
\end{ex}

\begin{ex} %[5D12V69]
Một người mua $2{,}5\text{ kg}$ thịt lợn giá $120\,000\text{ đồng/kg}$ và $3\text{ kg}$ cá tươi giá $80\,000\text{ đồng/kg}$. Người đó đưa cho người bán hàng $600\,000\text{ đồng}$. Số tiền người bán hàng phải trả lại là
\choice
{$50\,000\text{ đồng}$}
{\True $60\,000\text{ đồng}$}
{$70\,000\text{ đồng}$}
{$80\,000\text{ đồng}$}
\loigiai{
Số tiền mua thịt lợn là:
$$120\,000 \times 2{,}5 = 300\,000\text{ (đồng)}.$$
Số tiền mua cá tươi là:
$$80\,000 \times 3 = 240\,000\text{ (đồng)}.$$
Tổng số tiền phải trả là:
$$300\,000 + 240\,000 = 540\,000\text{ (đồng)}.$$
Số tiền được trả lại là:
$$600\,000 - 540\,000 = 60\,000\text{ (đồng)}.$$
\textbf{Đáp án của câu hỏi là B}
}
\end{ex}

\begin{ex} %[5D12V69]
Tính nhanh: $\dfrac{1}{2 \times 3} + \dfrac{1}{3 \times 4} + \dfrac{1}{4 \times 5} + \dots + \dfrac{1}{9 \times 10}$. Kết quả là
\choice
{$\dfrac{3}{10}$}
{\True $\dfrac{2}{5}$}
{$\dfrac{9}{20}$}
{$\dfrac{1}{2}$}
\loigiai{
Ta có:
$$A = \left(\dfrac{1}{2} - \dfrac{1}{3}\right) + \left(\dfrac{1}{3} - \dfrac{1}{4}\right) + \dots + \left(\dfrac{1}{9} - \dfrac{1}{10}\right) = \dfrac{1}{2} - \dfrac{1}{10} = \dfrac{4}{10} = \dfrac{2}{5}.$$
\textbf{Đáp án của câu hỏi là B}
}
\end{ex}

\begin{ex} %[5D12V69]
Kết quả của phép tính $12{,}5 \times 0{,}8 \times 2{,}5 \times 4$ là
\choice
{$80$}
{\True $100$}
{$120$}
{$150$}
\loigiai{
Nhóm các thừa số thích hợp:
$$(12{,}5 \times 0{,}8) \times (2{,}5 \times 4) = 10 \times 10 = 100.$$
\textbf{Đáp án của câu hỏi là B}
}
\end{ex}

\begin{ex} %[5D12V69]
Tìm $y$ biết: $y : 0{,}25 + y : 0{,}5 = 36$.
\choice
{$5$}
{\True $6$}
{$8$}
{$4{,}5$}
\loigiai{
Ta có $y : 0{,}25 = y \times 4$ và $y : 0{,}5 = y \times 2$.\\
Biểu thức trở thành:
$$y \times 4 + y \times 2 = 36$$
$$y \times 6 = 36 \Rightarrow y = 6.$$
\textbf{Đáp án của câu hỏi là B}
}
\end{ex}

\begin{ex} %[5D12V69]
Một hình chữ nhật có chu vi $28\text{ m}$. Nếu bớt chiều dài đi $2\text{ m}$ và thêm vào chiều rộng $2\text{ m}$ thì trở thành hình vuông. Diện tích hình chữ nhật ban đầu là
\choice
{$45\text{ m}^2$}
{\True $45\text{ m}^2$}
{$48\text{ m}^2$}
{$40\text{ m}^2$}
\loigiai{
Nửa chu vi hình chữ nhật là:
$$28 : 2 = 14\text{ (m)}.$$
Bớt chiều dài $2\text{ m}$, thêm chiều rộng $2\text{ m}$ thì hai cạnh bằng nhau, nghĩa là chiều dài hơn chiều rộng là:
$$2 + 2 = 4\text{ (m)}.$$
Chiều dài ban đầu là:
$$(14 + 4) : 2 = 9\text{ (m)}.$$
Chiều rộng ban đầu là:
$$14 - 9 = 5\text{ (m)}.$$
Diện tích hình chữ nhật ban đầu là:
$$9 \times 5 = 45\text{ (m}^2\text{)}.$$
\textbf{Đáp án của câu hỏi là B}
}
\end{ex}

\begin{ex} %[5D12V69]
Thực hiện phép tính: $\left(1 - \dfrac{1}{2}\right) \times \left(1 - \dfrac{1}{3}\right) \times \left(1 - \dfrac{1}{4}\right) \times \dots \times \left(1 - \dfrac{1}{100}\right)$. Kết quả là
\choice
{$\dfrac{1}{50}$}
{\True $\dfrac{1}{100}$}
{$\dfrac{2}{99}$}
{$\dfrac{99}{100}$}
\loigiai{
Biểu thức trở thành:
$$\dfrac{1}{2} \times \dfrac{2}{3} \times \dfrac{3}{4} \times \dots \times \dfrac{99}{100} = \dfrac{1}{100}.$$
\textbf{Đáp án của câu hỏi là B}
}
\end{ex}

\begin{ex} %[5D12V69]
Hiệu của hai số là $18{,}6$. Nếu gấp số lớn lên $3$ lần và giữ nguyên số bé thì hiệu mới là $78{,}6$. Số lớn ban đầu là
\choice
{$28$}
{\True $30$}
{$32$}
{$35$}
\loigiai{
Khi gấp số lớn lên $3$ lần thì hiệu tăng thêm $2$ lần số lớn.\\
Hai lần số lớn là:
$$78{,}6 - 18{,}6 = 60.$$
Số lớn ban đầu là:
$$60 : 2 = 30.$$
\textbf{Đáp án của câu hỏi là B}
}
\end{ex}

\begin{ex} %[5D12V69]
Biết rằng $0{,}75\text{ lít}$ sữa cân nặng $0{,}78\text{ kg}$. Hỏi $4{,}5\text{ lít}$ sữa như thế cân nặng bao nhiêu ki-lô-gam?
\choice
{$4{,}58\text{ kg}$}
{\True $4{,}68\text{ kg}$}
{$4{,}72\text{ kg}$}
{$4{,}8\text{ kg}$}
\loigiai{
Số lần $4{,}5\text{ lít}$ gấp $0{,}75\text{ lít}$ là:
$$4{,}5 : 0{,}75 = 6\text{ (lần)}.$$
Khối lượng của $4{,}5\text{ lít}$ sữa là:
$$0{,}78 \times 6 = 4{,}68\text{ (kg)}.$$
\textbf{Đáp án của câu hỏi là B}
}
\end{ex}

\begin{ex} %[5D12V69]
Giá trị của biểu thức $(15{,}8 + 24{,}2) : 5 - 2{,}4 \times 1{,}5$ là
\choice
{$4{,}2$}
{\True $4{,}4$}
{$4{,}6$}
{$5{,}0$}
\loigiai{
Ta có: $15{,}8 + 24{,}2 = 40$.\\
$40 : 5 = 8$.\\
$2{,}4 \times 1{,}5 = 3{,}6$.\\
Giá trị biểu thức là:
$$8 - 3{,}6 = 4{,}4.$$
\textbf{Đáp án của câu hỏi là B}
}
\end{ex}

\textbf{Phần 2. Trắc nghiệm đúng sai}
\setcounterpageref{ex}{0}

\begin{ex} %[5D12V69]
Cho biểu thức $M = 3{,}7 \times 48 + 3{,}7 \times 53 - 3{,}7$. Xét tính đúng/sai của các nhận định sau:
\choiceTF
{\True Biểu thức có thể đặt thừa số chung là $3{,}7$}
{\True Biểu thức tương đương với $3{,}7 \times (48 + 53 - 1)$}
{Giá trị của biểu thức $M$ bằng $37$}
{\True Giá trị của biểu thức $M$ bằng $370$}
\loigiai{
\begin{itemchoice}
\item Đúng. Có thừa số $3{,}7$ chung trong tất cả các số hạng.
\item Đúng. $M = 3{,}7 \times (48 + 53 - 1)$.
\item Sai. Vì $48 + 53 - 1 = 100 \Rightarrow 3{,}7 \times 100 = 370$ (chứ không phải $37$).
\item Đúng. $M = 370$.
\end{itemchoice}
}
\end{ex}

\begin{ex} %[5D12V69]
Xét tính đúng/sai của các kết quả phép tính sau:
\choiceTF
{\True $12{,}5 \times 0{,}4 = 5$}
{\True $4{,}8 : 0{,}2 = 24$}
{$\dfrac{3}{5} \times \dfrac{5}{7} + \dfrac{2}{7} = 1$}
{\True $2{,}5 \times 4 - 3{,}6 : 0{,}6 = 4$}
\loigiai{
\begin{itemchoice}
\item Đúng. $12{,}5 \times 0{,}4 = 5$.
\item Đúng. $4{,}8 : 0{,}2 = 4{,}8 \times 5 = 24$.
\item Sai. $\dfrac{3}{5} \times \dfrac{5}{7} = \dfrac{3}{7} \Rightarrow \dfrac{3}{7} + \dfrac{2}{7} = \dfrac{5}{7} \neq 1$.
\item Đúng. $2{,}5 \times 4 = 10$; $3{,}6 : 0{,}6 = 6 \Rightarrow 10 - 6 = 4$.
\end{itemchoice}
}
\end{ex}

\textbf{Phần 3. Câu trả lời ngắn}
\setcounterpageref{ex}{0}

\begin{ex} %[5D12V69]
Tính giá trị của biểu thức: $125 \times 17 \times 8$.
\shortans{17000}
\loigiai{
Áp dụng tính chất giao hoán và kết hợp:
$$125 \times 17 \times 8 = (125 \times 8) \times 17 = 1\,000 \times 17 = 17\,000.$$
}
\end{ex}

\begin{ex} %[5D12V69]
Tìm $x$ biết: $(x - 15) \times 4 = 160$.
\shortans{55}
\loigiai{
Ta có:
$$x - 15 = 160 : 4$$
$$x - 15 = 40 \Rightarrow x = 40 + 15 = 55.$$
}
\end{ex}

\begin{ex} %[5D12V69]
Một người nuôi $48$ con thỏ. Người đó đã bán đi $\dfrac{3}{8}$ số thỏ đó. Hỏi người đó còn lại bao nhiêu con thỏ?
\shortans{30}
\loigiai{
Số thỏ đã bán là:
$$48 \times \dfrac{3}{8} = 18\text{ (con)}.$$
Số thỏ còn lại là:
$$48 - 18 = 30\text{ (con)}.$$
}
\end{ex}

\begin{ex} %[5D12V69]
Tính giá trị của biểu thức: $1{,}2 \times 45 + 1{,}2 \times 55$.
\shortans{120}
\loigiai{
Đặt thừa số chung:
$$1{,}2 \times (45 + 55) = 1{,}2 \times 100 = 120.$$
}
\end{ex}

\begin{ex} %[5D12V69]
Một mảnh đất hình chữ nhật có chiều dài $24\text{ m}$, chiều rộng bằng $\dfrac{2}{3}$ chiều dài. Diện tích của mảnh đất đó là bao nhiêu mét vuông?
\shortans{384}
\loigiai{
Chiều rộng của mảnh đất là:
$$24 \times \dfrac{2}{3} = 16\text{ (m)}.$$
Diện tích của mảnh đất là:
$$24 \times 16 = 384\text{ (m}^2\text{)}.$$
}
\end{ex}

\begin{ex} %[5D12V69]
Cho hai số có tổng bằng $240$. Biết $\dfrac{1}{3}$ số thứ nhất bằng $\dfrac{1}{5}$ số thứ hai. Số thứ hai là bao nhiêu?
\shortans{150}
\loigiai{
Tỉ số giữa số thứ nhất và số thứ hai là $\dfrac{3}{5}$.\\
Tổng số phần bằng nhau là:
$$3 + 5 = 8\text{ (phần)}.$$
Số thứ hai là:
$$240 : 8 \times 5 = 150.$$
}
\end{ex}
